% Page 070 — page imprimée 66
% La famille Roussel du Villaret au Château d’Ayres (suite)

\textbf{Il était une fois au Villaret}, hameau de la vallée du Bétuzon à sept km de Meyrueis une
famille de paysans, les Roussel qui possédaient une petite maison avec une grange attenante
actuellement rénovée par les Danjean-Henry. La vie était très dure mais on économisait sou
à sou et arriva le moment où la famille envisagea de descendre « à la ville » pour avoir une
vie plus facile. La ville c’était Meyrueis, bien sûr !

\medskip

Comment purent-ils acheter un café au bord de leur rivière, le Bétuzon ? Et à quelle époque ?
L’histoire ne le dit pas. Mais Antoine Roussel avait épousé une Valès, qui avait quelques
biens, sœur d’un certain Henry Valès, né à Meyrueis qui fit de très brillantes études de
Théologie à Montauban, d’où il sortit « major ». Il épousa en 1861 « un beau parti », Suzanne
Martin, fille de propriétaire de Meyrueis qui lui apporta 20.000 F en dot. Nous ignorons
pourquoi il s’expatria aux Pays Bas où il fut nommé pasteur de l’Eglise Wallonne et même,
disait-on, chapelain de la famille royale (Les Orange-Nassau) alors que Wilhelmine était
encore toute jeune. Il repose, au cimetière de Meyrueis, dans un caveau surmonté d’une belle
colonne brisée entourée de lierre et sculptée par des ouvriers venus tout spécialement de
Hollande !!

\medskip

Or, cet Henry Valès était propriétaire du « Café National » surplombant le Bétuzon derrière la
halle. Les Roussel ont-ils hérité en partie de ce café ? C’est probable ! Antoine a-t-il racheté la
part de son beau frère ? En 1880, 1900 ? Nous ne le savons pas encore.

\medskip

Aux alentours de 1900 se réunissaient dans ce café les républicains de Meyrueis. Antoine
Roussel, propriétaire et débitant se déplaçait même en Italie pour acheter du vin de qualité
sans doute à un prix compétitif !

\medskip

Son fils, Jules, « monta » à Paris pour poursuivre des études. Il devint représentant de la
maison Lauret, gantiers de Millau qui avaient une excellente réputation. Il avait donc ses
entrées dans de grandes maisons de couture, Hermès par exemple, et c’est sans doute grâce à
de multiples connaissances qu’il fit la rencontre d’une demoiselle Hélène Harrison. Elle était
de « bonne famille », disposait d’une certaine fortune. Fut-ce un coup de foudre ? Nous
l’ignorons mais, toujours est-il qu’ils s’épousèrent et que le jeune couple s’installa à Meyrueis
pour les vacances.

\medskip

Arrivant de Paris la jeune femme trouva que la maison, l’ancien café, était bien étriquée, peut-
être en piètre état. Probablement sans eau ni électricité, (elle fut installée en 1920 seulement).
Il se trouva que le château d’Ayres était en vente, il séduisit Hélène Roussel : réunissant leurs
liquidités, le jeune couple put en faire l’achat.

\medskip

\textbf{La restauration du château d’Ayres}

\medskip

Cette grande bâtisse n’ayant guère été entretenue par ses précédents propriétaires, la famille
Nogaret, Jules Roussel s’employa aussitôt à la remettre en état, payant de sa personne avec
les maçons de Meyrueis. Les salons, la salle à manger, et les nombreuses chambres purent
accueillir des hôtes – payants (il fallait bien vivre) et rapidement le château d’Ayres devint
une pension-hôtel renommée comme on disait à l’époque, entre les deux guerres, avec
piscine (un grand bassin rectangulaire et arrondi aux deux extrémités devant le château avec
son jet d’eau central), tennis, tout cela dans un parc de 5ha agrémenté de magnifiques
châtaigniers et de quelques séquoias imposants.
